\section{Experimental method}
\label{sec:method}

% Slides 1.3 and 2.9: define objectives and hypotheses, the setting, the
% measure, then analyse with statistics and plots.

\subsection{Setting}
% Maps: tuned on l_hall, evaluated on dead_end and l_turn, which no model
% was tuned on (avoids the tuning-map confound). Seeds 0-9: each seed is
% one spawn, shared by every model and setting, so runs pair by seed.
% Two experiments: (1) the evaluation, n = 20, j = 2, model x plan on both
% held-out maps; (2) the sweep, dead_end with the plan, n x j x model.
\todo{Maps, seeds, the two experiments.}

\subsection{Measure}
% Success: 90% of the drones in C within the limit, go time + twice the
% leaders' own time along their planned path (83 s on dead_end, 112 s on
% l_turn). Why a limit: a flock that lost its leaders wanders into C by
% chance given time (false positives, design log 2026-10-28). Why twice:
% larger flocks pass gaps slower. Also logged: time to 90%, the closest
% two drones came, crashes.
\todo{Success within the limit; secondary measures.}

\subsection{Statistics}
% Paired by seed. Success is binary: exact McNemar test (binomial on the
% discordant seeds). Closest pair: Wilcoxon signed-rank. Leader split:
% Fisher's exact test. Uncorrected p, many comparisons: read as evidence,
% not proof. With 10 seeds the smallest McNemar p is 0.002.
\todo{Tests and their limits.}

\subsection{Traps, and how they were handled}
% Nehmzow's traps (slides 1.3):
%   confounds -> held-out maps; models tuned by the same protocol;
%   floor/ceiling -> l_turn is at the ceiling, dead_end separates; the
%     j = 0 control fixes the floor;
%   systematic errors -> identical spawns per seed across models;
%   wishful thinking -> hypotheses before the sweep, the limit kept when
%     it cost results ("Too hard a limit? Checked, kept").
\todo{Traps.}
